\section{関連研究}
\label{sec:related-work}

\subsection{ジオソーシャルデータと地理空間分析}

位置情報を伴う投稿や地物データは，人間の活動と都市・自然環境の関係を分析するための情報源として利用されている．写真共有サービスを対象とした研究では，KisilevichらがFlickrとPanoramioのジオタグ付き写真からイベントや人間行動を分析し\cite{kisilevich2010eventbased}，K{\'a}d{\'a}rとGedeが写真の位置分布に基づいて観光客の空間行動を可視化した\cite{kadar2013wheredotouristsgo}．Huらはジオタグ付き写真から都市のArea of Interest (AOI) を抽出し，タグや写真を用いてその特徴を分析している\cite{hu2015extractingAOI}．著者らも，撮影地点分布から撮影スポットとAOIを検出する手法\cite{shirai2013aoi}や，Flickr上のタグに対応する位置分布から海岸線を抽出する手法\cite{omori2014coastlines}を検討してきた．

地理空間分析の対象は写真投稿に限られない．例えば，OpenStreetMap (OSM) を利用するOSMnxは，道路網や建物などの地理データの取得・構築・分析・可視化を支援している\cite{boeing2017osmnx}．これらの研究を踏まえ，本研究では，写真投稿に加えて施設などのPoint of Interest (POI) も扱い，カテゴリごとの空間分布を探索できるソフトウェアを対象とする．写真の投稿地点と施設の所在地は異なる現象を表すため，データ形式や可視化手順を共通化することと，観測の意味を同一視することは区別する．

\subsection{外部APIを介したデータ取得とサービス固有の制約}

外部サービスから地理データを取得するソフトウェアには，取得先に応じた問い合わせの生成と，応答を分析可能な形式へ変換する役割がある．Foxらの\textit{photosearcher}はFlickrの写真とメタデータをRから取得するためのパッケージであり，検索条件の指定や大規模な検索結果の取得を支援する\cite{fox2020photosearcher}．Padghamらの\textit{osmdata}はOverpass APIへの問い合わせによりOSMデータを取得し，Rの\textit{sf}または\textit{sp}で処理可能な空間オブジェクトへ変換する\cite{padgham2017osmdata}．また，OSMnxはデータ取得から道路ネットワークの構築・分析までを一貫した処理として提供する\cite{boeing2017osmnx}．これらは，サービス固有のアクセス方法と後段の分析処理を接続する，関連するソフトウェア実装である．

取得戦略を検討する具体例として，Flickrの検索APIにはページングに関する制約が報告されている．Foxらは，一つの検索条件で取得できるユニークな写真が最大4,000件に制限され，それを超えるページで既取得の写真が重複して返される問題を指摘した\cite{fox2020photosearcher}．Yanらは検索領域を再帰的に細分化してこの制約に対処している\cite{yan2017postdisaster,yan2019flickrtags}．一方，\textit{photosearcher}は時刻順に取得した最大4,000件の写真の最新撮影日時を次の検索下限へ設定し，検索範囲を時間方向へ進める\cite{fox2020photosearcher}．したがって，空間分割や取得済みデータの時刻に基づく検索境界の更新自体は，既存研究にも見られる考え方である．

さらに，Gulnermanらは同じ期間に対するFlickrの繰り返し検索で結果が変化することを報告し\cite{gulnerman2024antarctica}，Lepp{\"a}m{\"a}kiらは同一の時空間条件に対する応答の不一致と，検索領域の細分化による取得結果の安定性の改善を示している\cite{leppamaki2025flickr}．これらの知見は，検索結果件数やページ番号だけを前提とした単一の取得方式を，異なるサービスへ一律に適用することの限界を示唆する．本研究では，Flickrへの対処を基盤全体の前提とするのではなく，取得先に応じて収集戦略を実装できる構成を重視する．

\subsection{地理空間処理と可視化の再利用}

取得後の処理では，データを共通の地理空間表現へ変換することが処理の再利用を支える．Pebesmaの\textit{sf}は，点・線・ポリゴンとその属性を標準的なSimple Featuresとして扱い，入出力，幾何演算，集約，可視化などをRの分析環境へ接続する\cite{pebesma2018sf}．このような共通表現の利用は，データ取得先と分析処理を分離するための基礎となる．

また，前述の写真分布の可視化\cite{kisilevich2010eventbased,kadar2013wheredotouristsgo}や，OSMnxによる地理データの可視化\cite{boeing2017osmnx}は，取得結果を探索可能な表現へつなぐことの重要性を示している．本研究が目指すのは，こうした個々の分析手法や汎用的な地理情報処理ライブラリの置き換えではない．対象領域とカテゴリの指定，サービスへの問い合わせ，取得結果の空間集約，複数の地図表現による閲覧を，一つの操作系で利用できるようにすることである．

\subsection{Splatoneの位置づけ}

Splatoneは，サービス別のデータ取得を担うProviderと，取得データの地図表現を担うVisualizerを独立したプラグインとして構成する．Flickr，Google Places Text Search，OSMのOverpass APIに対応するProviderは，それぞれの検索条件や応答に応じた取得処理を行い，共通のGeoJSONベースの表現へ変換する．後段では，対象領域，カテゴリ，Hexグリッドを共有しながら，点表示，クラスタ表示，ヒートマップ，カテゴリ別のセル集約などのVisualizerを利用する．これにより，新しい取得先への対応と，新しい可視化方法への対応を分離する．

Flickr Providerは，この分離の意義を示す具体例である．Hexセルを包含するbounding boxを空間検索に用い，時刻降順の第1ページのみを取得して，取得結果の時刻に基づき次の検索上限を更新する．この\textit{pagination-independent}な方式は，後続ページを走査せずに広域の収集を進めるための，Flickr固有の実装上の工夫として位置づける．一方，Google PlacesとOverpassのProviderには，それぞれのAPIに適した取得処理を実装し，Flickrの時間カーソル方式を要求しない．

したがって，本研究の中心的な貢献は，時間に基づく検索やプラグイン化という考え方そのものではなく，異種APIに対応する取得戦略と，共通のカテゴリ・空間表現および交換可能な可視化を結び付けた，拡張可能な収集・可視化ソフトウェアを提供する点にある．Flickr Providerでは局所的な同時刻の大量投稿を完全には列挙しない場合があり，広い時空間範囲で収集を継続することを優先する．この設計判断はProvider固有の制約として明示し，収集結果の完全性や，省略が分析へ与える影響の小ささを，基盤全体の保証とはしない．

% Software description / Flickr Provider へ統合するための実装メモ：
% - 各Hexの外接BBOXは重複するが，取得後に元のHex内の点だけを採用する．
%   BBOXの重複によって，同一時刻の取りこぼしが回収されるとは主張しない．
% - per_page=250, page=1．DateModeに応じてdate-posted-desc/date-taken-desc．
% - 次のDateMaxは取得結果の最古時刻．全件同一時刻なら基本更新は1秒後退．
% - 単一ownerによる短時間の大量投稿には別途追加の時間スキップがある．
%   これは同一時刻の1秒後退と区別して実装章で説明する．
% - Haste有効時には残り時間範囲を分割する．
% - 2014年の約2億枚の収集経験は現行Splatoneの性能評価と同一視しない．
%   Omori et al. (2014)を，論文中で報告されていない総収集数の根拠にしない．
% - 未測定の取りこぼし率・スパム率・分析への影響を定量的な結論にしない．
